% SemEval-2027 conditionally accepted task proposal: expanded revision.
% Revised proposal for SoftConf upload.

\documentclass[11pt]{article}
\usepackage[preprint]{acl}
\usepackage{times}
\usepackage{latexsym}
\usepackage[T1]{fontenc}
\usepackage[utf8]{inputenc}
\usepackage{microtype}
\usepackage{inconsolata}
\usepackage{booktabs}
\usepackage{amsmath}
\usepackage{graphicx}
\usepackage{url}

\newcommand{\EquivLabel}{\textsc{Equiv}}
\newcommand{\ForwardLabel}{\textsc{Forward}}
\newcommand{\BackwardLabel}{\textsc{Backward}}
\newcommand{\OtherLabel}{\textsc{Other}}
\newcommand{\SoftCons}{\textsc{SoftCons}}
\newcommand{\HardCons}{\textsc{HardCons}}
\newcommand{\Rev}{\mathrm{Rev}}

\title{SemEval-2027 Task Proposal: DiCo-NLI---Directional Consistency in Fine-Grained Natural Language Inference}
\author{
  Jon F. Apaolaza Larraya \and Aitor Soroa \and Rodrigo Agerri \and Inigo Lopez-Gazpio \\
  HiTZ Basque Center for Language Technology--Ixa NLP Group \\
  University of the Basque Country UPV/EHU \\
  \texttt{\{jonfelix.apaolaza, a.soroa, rodrigo.agerri, inigo.lopez\}@ehu.eus}
}

\begin{document}
\maketitle

\section{Overview}

\paragraph{Task.}
DiCo-NLI evaluates direction-sensitive fine-grained natural language inference (NLI) over ordered phrase pairs.
For each source pair $(A,B)$, systems predict a label for both the original order $(A,B)$ and the reversed order $(B,A)$.
The reversible label algebra is deterministic: \EquivLabel{} maps to itself, while \ForwardLabel{} and \BackwardLabel{} map to each other.
The shared task therefore measures both item-level label correctness and whether predictions remain compatible under reversal.

\paragraph{Motivation.}
NLI systems and large language models can obtain strong aggregate scores while relying on shallow cues, annotation artifacts, or incomplete semantic representations~\cite{gururangan2018annotation,lopez2024revisiting}.
Direction is a central part of NLI meaning: if \emph{everyone is hungry} entails \emph{someone is hungry}, the reversed ordered pair does not have the same label.
DiCo-NLI targets this direction-sensitive failure mode in a compact, auditable, and multilingual setting.

The task is inspired by reversal-based evaluation, including the Reversal Curse of \citet{berglund2024reversal}, but it is not a direct benchmark of parametric factual reversal.
Berglund et al.\ evaluate whether a model trained on a fact such as ``A is B'' can retrieve or generate the reverse form ``B is A'' from its parameters.
DiCo-NLI instead supplies both phrases and asks for their ordered inferential relation.
The connection is methodological rather than identical: reversal is used as a controlled stress test for directional generalization.
Follow-up work has studied reversal failures and mitigation strategies in factual recall and question-answering settings~\cite{golovneva2024reverse,lin2024delving}.
DiCo-NLI complements that line by operationalizing reversal in supervised fine-grained NLI.

\paragraph{Revision focus.}
This revised proposal addresses the conditional-acceptance reviews by sharpening the conceptual scope, renaming the task from CoCo-NLI to DiCo-NLI, replacing ``coherence'' with the more precise term ``consistency'', making the metric roles explicit, adding pilot results, specifying contamination controls, defining multilingual tracks, and detailing translation validation and baseline resources.
The official task target is directional consistency in fine-grained NLI, not a claim that DiCo-NLI scores are interchangeable with Berglund-style parametric Reversal Curse scores.

\paragraph{Expected participants and impact.}
The task should attract researchers working on NLI, semantic evaluation, multilingual NLP, interpretability, model robustness, and LLM reasoning diagnostics.
It contributes a reusable benchmark where label correctness and reversal consistency can be analyzed separately, with English, Spanish, Basque, and mixed-language tracks, an official scorer, validation tools, and runnable starter baselines.
The task is designed to remain useful after SemEval as a diagnostic resource for directional generalization and multilingual fine-grained inference.

\section{Relation to Prior Work}

\paragraph{Fine-grained NLI and PhrasIS.}
DiCo-NLI builds on Natural Logic-style fine-grained inference labels~\cite{maccartney2009natural} and on PhrasIS, a phrase inference and similarity benchmark derived from naturally occurring image captions and news headlines~\cite{lopezgazpio2024phrasis}.
PhrasIS reannotates phrase pairs from the Interpretable STS lineage~\cite{lopezgazpio2017interpretable}, where sentence pairs are decomposed into aligned semantic chunks such as nominal expressions, verb chains, prepositional phrases, adverbial expressions, and quantificational expressions.
DiCo-NLI is compositional in this phrase-level sense: it evaluates inference over local semantic units built from smaller lexical and syntactic material.
It is not intended as a controlled benchmark of arbitrary compositional generalization over generated templates.

\paragraph{Consistency-based NLI evaluation.}
DiCo-NLI is also related to consistency evaluation for language models.
BECEL~\cite{jang2022becel} evaluates several behavioral consistency properties, including symmetric consistency for standard three-way NLI.
DiCo-NLI differs by explicitly modeling direction-sensitive label reversal in a fine-grained label space: \texttt{FORWARD\_ENTAILMENT} and \texttt{BACKWARD\_ENTAILMENT} swap under reversal, while \texttt{EQUIVALENCE} is invariant.
Recent work on NLI meta-inferential properties~\cite{blanck2026reverse} and transitive self-consistency~\cite{wu2025transitive} is complementary, but remains in the research-benchmark setting.
DiCo-NLI brings a multilingual, fine-grained, directional-reversal setup into a shared-task format with official data, scorer, validation tools, and participant submissions.

\section{Data and Tracks}

\subsection{Source Data and Label Algebra}

The task starts from the reversible subset of PhrasIS.
The English source data contains 1,946 human-annotated source pairs and 3,892 ordered instances after materializing both directions.
The three reversible labels are close to balanced (Table~\ref{tab:label-data}), reducing the risk that paired correctness can be inflated by a trivial majority-class strategy.

\begin{table}[t]
\centering
\small
\begin{tabular}{@{}lrrl@{}}
\toprule
Label & Source & Ordered & $\Rev(\cdot)$ \\
\midrule
\texttt{EQUIVALENCE} & 686 & 1,372 & \texttt{EQUIVALENCE} \\
\texttt{FORWARD} & 624 & 1,260 & \texttt{BACKWARD} \\
\texttt{BACKWARD} & 636 & 1,260 & \texttt{FORWARD} \\
\midrule
\textbf{Total} & \textbf{1,946} & \textbf{3,892} & -- \\
\bottomrule
\end{tabular}
\caption{English reversible source data and label-reversal algebra. ``Ordered'' counts both directions.}
\label{tab:label-data}
\end{table}

The official label inventory contains four code values:
\texttt{EQUIVALENCE}, \texttt{FORWARD\_ENTAILMENT}, \texttt{BACKWARD\_ENTAILMENT}, and \texttt{NEGATIVE\_OTHER}.
The first three labels form the reversible algebra.
Original PhrasIS labels outside this algebra will be pooled as \texttt{NEGATIVE\_OTHER}.
These items make the classification problem less artificially clean and help systems distinguish directional entailment from weaker or unrelated semantic relations.
They enter item-level label evaluation but are excluded from \SoftCons{} and \HardCons{}, which are defined only over reversible source pairs.
\texttt{NEGATIVE\_OTHER} items will be capped and stratified so that they do not dominate the reversible labels; final counts and per-split distributions will be documented for every track.

\subsection{Tracks and Scale}

DiCo-NLI will include four independent submission tracks (Table~\ref{tab:tracks}).
Participants may enter any subset of tracks.
The three monolingual tracks evaluate standard within-language phrase-pair NLI.
Cross-lingual transfer is handled explicitly in the Mixed track, where premise and hypothesis may use different languages.
This design avoids ambiguity about whether EN--EU transfer is hidden inside the monolingual tracks and makes macro-averages interpretable.

\begin{table}[t]
\centering
\small
\resizebox{\columnwidth}{!}{%
\begin{tabular}{@{}lll@{}}
\toprule
Track & Name & Description \\
\midrule
1 & English & Monolingual English phrase-pair NLI \\
2 & Spanish & Monolingual Spanish phrase-pair NLI \\
3 & Basque & Monolingual Basque phrase-pair NLI \\
4 & Mixed & Premise and hypothesis use different EN/ES/EU languages \\
\bottomrule
\end{tabular}
}
\caption{Planned DiCo-NLI tracks.}
\label{tab:tracks}
\end{table}

Before quality filtering, the reversible component yields up to 11,676 monolingual ordered instances across EN/ES/EU and 23,352 mixed-language ordered instances, for 35,028 reversible ordered instances across all tracks.
Final counts will be lower if translation validation removes unstable target-language pairs.
If filtering distorts a target-language label distribution, controlled stratified downsampling will restore comparable proportions and the adjustment will be reported.

\subsection{Splits, Leakage, and Contamination}

The official split unit is the underlying source phrase pair, not the ordered instance.
For every source pair $(A,B)$, all derived records are assigned to the same train, development, trial, or evaluation partition: the original direction, the reversed direction, Spanish and Basque translations, and every mixed-language variant.
Thus, if a source pair appears in the final evaluation split, no direction, language, or variant of that pair appears in released training or development data.
This implements the hard/unseen setting from the LREC pilot as the official SemEval design.

We do not claim that DiCo-NLI is contamination-free.
PhrasIS has been publicly described since 2024, and frontier models may have seen some examples or documentation.
The task nevertheless reduces avoidable leakage by using new grouped splits, hidden final labels during evaluation, delayed gold release, complete provenance metadata, newly produced multilingual variants, and exact/near-exact overlap audits against examples appearing in papers, public documentation, and released sample material.
The official data release will include a \texttt{data/README.md} documenting split construction, source-pair grouping, label distributions, provenance, and contamination-control decisions.

\subsection{Multilingual Production and Validation}

The multilingual tracks require more than phrase-level translation fidelity: the translated pair must preserve the intended NLI relation and its deterministic reverse relation.
Spanish and Basque candidates will be generated independently by multiple systems.
The initial generators are NLLB-200 and Google Cloud Translation, representing an open multilingual model and a production API.
For Basque, we will also use the most recent Latxa version available through a HiTZ API.
During internal preparation, we will additionally try selected open and closed translation-capable LLMs through OpenRouter or comparable providers when licensing, cost, and reproducibility constraints permit.
System/API versions, providers, prompts, decoding settings, and dates will be frozen and recorded.
Generators receive only the phrases, without relation labels or reverse metadata, so that translation is not biased toward the expected label.

Automatic agreement signals will be used only for triage.
High-disagreement cases will be inspected more carefully, but automatic overlap is not treated as evidence that an entailment relation is preserved.
For each target language, two bilingual annotators will independently assess a 100-source-pair pilot audit.
The audit will check: (i) phrase-level translation adequacy, (ii) preservation of the original NLI label, and (iii) preservation of the reverse label after swapping premise and hypothesis.
The audit will deliberately include, where available, examples involving quantification, plurality, definiteness, specificity, genericity, scope, case marking, clitics, and other morphosyntactically sensitive phenomena.
We will report raw agreement, Cohen's $\kappa$, adjudication counts, and translation-induced label-flip or instability rates.

For full production, one bilingual reviewer checks every translated item.
A second reviewer adjudicates low-confidence, ambiguous, or relation-changing cases.
Items are corrected when a relation-preserving translation can be found; otherwise they are removed from official evaluation.
This applies equally to apparently simple examples such as \emph{Everyone is hungry} / \emph{Someone is hungry}, because quantifier scope and generic/specific readings can be destabilized by translation.

\subsection{Release, Licenses, and Ethics}

The upstream PhrasIS repository is MIT-licensed; its notice and provenance will be retained.
Newly authored translations, adjudications, splits, and metadata will be released under CC BY 4.0, while scorer and starter-kit software will use GPL-3.0.
A component-level license manifest, GitHub repository, and permanent Zenodo record will accompany trial and final data.
Before release, final items will be screened for private-person personally identifiable information and sensitive content.
The task supports semantic classification research and involves neither personal profiling nor medical decision making.
The main methodological risk is benchmark overfitting to shallow cues, which is one of the behaviors the task is designed to expose.

\section{Pilot Evidence, Trial Data, and Baselines}

\subsection{Pilot Results}

The task is grounded in the public LREC 2026 pilot~\cite{apaolaza2026assessing}, which evaluated encoder-only, decoder-only, and encoder--decoder systems on English PhrasIS reversal pairs.
The pilot included DeBERTa-v3, ModernBERT, RoBERTa, OPT, Pythia, GPT-2, BART, and Flan-T5.
Table~\ref{tab:pilot} gives representative and family-level results.
Weighted $F_1$ is measured on the positive combined English test track; \SoftCons{} and \HardCons{} use the unseen/hard positive combined track, the condition adopted for SemEval.

\begin{table*}[t]
\centering
\small
\begin{tabular}{@{}llcc@{}}
\toprule
Family & Representative model & Representative $F/S/H$ & Family mean $\pm$ sd $F/S/H$ \\
\midrule
Encoder-only & DeBERTa-v3-base & .75/.84/.79 & $.713\!\pm\!.031$ / $.770\!\pm\!.049$ / $.721\!\pm\!.048$ \\
Decoder-only & OPT-125M & .61/.65/.58 & $.583\!\pm\!.038$ / $.568\!\pm\!.060$ / $.498\!\pm\!.070$ \\
Encoder--decoder & BART-large & .68/.80/.72 & $.655\!\pm\!.024$ / $.768\!\pm\!.025$ / $.693\!\pm\!.019$ \\
\bottomrule
\end{tabular}
\caption{Pilot results. $F$: weighted $F_1$; $S$: \SoftCons{}; $H$: \HardCons{}.}
\label{tab:pilot}
\end{table*}

These results demonstrate feasibility and leave room for improvement.
Encoder-only systems perform strongest overall, but even the best representative encoder does not solve the hard/unseen directional setting.
Decoder-only systems lag substantially, especially under paired correctness, and encoder--decoder systems show a different profile: their \SoftCons{} can be close to encoder-only models even when item-level $F_1$ is lower.
This supports reporting item-level and paired metrics together rather than collapsing evaluation to a single coarse score.

\subsection{Trial Data and Supplementary Materials}

For the revision package, we will provide pilot/trial data demonstrating the complete task format.
The trial release will be a source-pair grouped random sample from PhrasIS, including original and reversed records, labels, reverse identifiers, split metadata, source provenance, and language fields.
The sample will exercise all four labels in the official format.
If Spanish or Basque validation is not complete by the revision deadline, the trial data will clearly mark which records are English-only pilot records and which multilingual records are provisional examples.

The supplementary upload will also include baseline result sheets confirming that the evaluation code is operational.
The LREC pilot paper will be attached as the detailed baseline evidence for the broad model comparison, and the supplementary results will expose the compact table above in machine-readable form.
The evaluation plan will be provided as a separate document and will match the metric definitions in Section~\ref{sec:evaluation}.

\subsection{Starter Kit}

The starter kit is intended to provide runnable reference systems, not the strongest possible systems.
We will fix one reference per architecture family:
\path{microsoft/deberta-v3-base} for encoder-only classification, \path{facebook/opt-125m} for decoder-only classification or constrained generation, and \path{facebook/bart-base} for encoder--decoder generation/classification.
For each baseline, we will release preprocessing code, configuration files, prediction format, scorer command, learning rate, batch size, epochs, maximum sequence length, optimizer, random seeds, model-selection rule, and trained checkpoints when redistribution permits.
Data construction and official scoring are CPU-based; each starter baseline will use a documented one-GPU configuration.

\section{Evaluation}
\label{sec:evaluation}

\subsection{Official Scores}

Systems predict one label for every ordered phrase pair.
The official evaluation for each track reports three main scores: weighted $F_1$ over all four labels, \SoftCons{}, and \HardCons{}.
These scores are not collapsed into a single primary metric because they answer different questions.
Weighted $F_1$ measures item-level label correctness, \SoftCons{} measures directional self-consistency independently of correctness, and \HardCons{} measures paired directional correctness over reversible source pairs.
Macro-$F_1$ and per-label $F_1$ are also reported as secondary diagnostics.
Official results are track-specific.
We will additionally report a monolingual macro-average over EN/ES/EU and an all-track macro-average over EN/ES/EU/Mixed, but neither replaces the per-track leaderboards.

The scientific analysis will foreground the full weighted-$F_1$/\SoftCons{}/\HardCons{} profile: a system may be accurate but directionally unstable, internally stable but systematically wrong, or strong only when both directions of a reversible pair are considered together.
\HardCons{} is interpreted as a strict paired-correctness summary rather than as an independent consistency axis beyond correctness.
If the CodaBench interface requires a single display ordering or an additional summary field, that convention will be preannounced in the scorer documentation and will not replace the three official evaluation scores.

\subsection{Metric Definitions}

Let $L$ be the full label set, $n_\ell$ the support of label $\ell$, and $N_L=\sum_{\ell\in L}n_\ell$.
For each label $\ell$,
\[
P_\ell=\frac{TP_\ell}{TP_\ell+FP_\ell}, \qquad
R_\ell=\frac{TP_\ell}{TP_\ell+FN_\ell},
\]
\[
F_\ell=\frac{2P_\ell R_\ell}{P_\ell+R_\ell}, \qquad
F_w=\sum_{\ell\in L}\frac{n_\ell}{N_L}F_\ell .
\]
The scorer will follow standard zero-division handling and will document this behavior in the released evaluation plan.

Let $\mathcal{P}$ be the set of reversible source pairs in a track.
For pair $i\in\mathcal{P}$, let $x_i$ be the original ordered item, $x_i^r$ its reversed counterpart, $y_i$ the gold label for $x_i$, $\hat{y}_i=f(x_i)$ the system prediction, and $\hat{y}_i^r=f(x_i^r)$ the prediction for the reversed item.
The deterministic reversal operator $\Rev$ maps \texttt{EQUIVALENCE} to itself, \texttt{FORWARD\_ENTAILMENT} to \texttt{BACKWARD\_ENTAILMENT}, and \texttt{BACKWARD\_ENTAILMENT} to \texttt{FORWARD\_ENTAILMENT}.
\texttt{NEGATIVE\_OTHER} is not included in $\mathcal{P}$.

\begin{align*}
s_i &= \mathbf{1}[\hat{y}_i^r=\Rev(\hat{y}_i)],\\
\SoftCons{} &= |\mathcal{P}|^{-1}\sum_{i\in\mathcal{P}} s_i .
\end{align*}

\SoftCons{} measures directional compatibility independently of correctness.
For example, predicting \texttt{EQUIVALENCE} in both directions is soft-consistent even if the correct labels are \texttt{FORWARD\_ENTAILMENT} and \texttt{BACKWARD\_ENTAILMENT}.
Predicting \texttt{FORWARD\_ENTAILMENT} in both directions is not soft-consistent.

\begin{align*}
h_i &= \mathbf{1}[\hat{y}_i=y_i]\,
      \mathbf{1}[\hat{y}_i^r=\Rev(y_i)],\\
\HardCons{} &= |\mathcal{P}|^{-1}\sum_{i\in\mathcal{P}} h_i .
\end{align*}

\HardCons{} is paired directional correctness.
Because the reverse gold label is deterministic, correctness in both directions entails consistency automatically.
Thus, \HardCons{} is paired accuracy over reversal-defined semantic units.
It is still useful because it changes the unit of evaluation from isolated ordered instances to complete reversible source pairs: two systems can obtain the same item-level score while differing in whether they solve both directions of the same relation.

\subsection{Run Metadata and Rules}

Run metadata is mandatory.
Participants will report access regime (\texttt{open-weight}, \texttt{commercial/API}, \texttt{private/closed}, or \texttt{hybrid/ensemble}), architecture family, model/API identifier and version, model size when known, external data use, retrieval, prompting or fine-tuning strategy, and multilingual strategy.
Access regime supports post-task analysis and does not define a separate leaderboard.
Verification will be partial and based on system papers, released code/configurations, model/API documentation, and organizer sanity checks; unverifiable fields will be marked \texttt{undisclosed}.
Participants may use external resources unless the final rules restrict a specific resource, but all such use must be documented.

\section{Infrastructure, Schedule, and Organizers}

The task will run on CodaBench with a public repository containing trial data, scorer, format checker, starter kit, submission examples, and documentation.
Final gold labels will be hidden during the official evaluation phase and released afterward together with the scorer and data documentation.
Task artifacts will be archived on Zenodo for long-term access.

We will follow the SemEval-2027 schedule.
Sample data is due by 15 July 2026, training data by 1 September 2026, internal evaluation data by 1 December 2026, official evaluation in January 2027, system papers in February 2027, notifications in March 2027, camera-ready papers in April 2027, and the SemEval workshop in summer 2027.
The three-week revision period will be used to prepare the revised proposal, trial/pilot sample, baseline result sheets, evaluation plan, and supporting documentation in a private development repository before public release.

\paragraph{Organizers.}
\textbf{Inigo Lopez-Gazpio} (lead) works on semantic evaluation, STS, fine-grained semantics, and benchmark design, and co-organized SemEval-2015 Task~2, SemEval-2016 Task~2, and SemEval-2017 Task~1.
\textbf{Jon F. Apaolaza Larraya} (co-organizer) led the DiCo-NLI pilot and will coordinate data, scoring, and analysis.
\textbf{Aitor Soroa} (advisory) contributes expertise in semantic processing and robust evaluation.
\textbf{Rodrigo Agerri} (advisory) contributes multilingual NLP and system evaluation expertise.
All organizers are affiliated with the HiTZ Basque Center for Language Technology--Ixa NLP Group, University of the Basque Country UPV/EHU.
The team will maintain participant communication and infrastructure, release data and software, manage submissions and reviewing, and write the task description paper.

\bibliography{dico_nli_biblio}

\end{document}
